\section{Discussion}
\label{sec:discussion}

\paragraph{The topology as mechanistic account, not just a heuristic.} The forgetting map is more than a tuning trick. An \fvs{} distribution that is bimodal, with an early-layer / anchor-head resistance region and a late-MLP vulnerability region, is consistent with a reading in which early layers encode task-agnostic linguistic competence while late-layer MLPs act as task-specific key-value stores \citep{geva2023dissecting}. Fine-tuning that writes into the latter over-writes prior tasks' keys; restricting updates to the former avoids this. Reading \fvs{} as an \textit{explanation} rather than a selection rule also recasts catastrophic forgetting as a question of \textit{where in the network representational real-estate is being re-used}, which is a more tractable question than the global ``why does the model forget'' framing.

\paragraph{Causal vs.\ gradient attribution.} Our results arbitrate between two attribution families that have so far been compared informally. Causal-intervention \fvs{} identifies \textit{which} component's update causally degrades prior-task accuracy under counterfactual reversion. Gradient-attribution \fvs$_{\mathrm{grad}}$ identifies components whose per-example gradients have high magnitude on prior-task loss. These differ when a component's updates have small gradient but high downstream leverage, or vice versa. Our Spearman correlation between the two (\pend{causal_vs_grad_spearman}) quantifies how much they agree in practice.

\paragraph{Connection to the EMNLP 2026 explainability theme.} Our work exemplifies a pragmatic form of explainability: not a post-hoc rationalisation of an individual prediction, but a structural account of \textit{where} capability lives and thus \textit{where} it can safely be updated. This reframes PEFT target selection as an interpretability-grounded design choice rather than a convention inherited from the original LoRA paper.

\paragraph{Failure modes.} When task-relevant knowledge is stored in the forgetting-vulnerable region, \tglora{} under-performs standard LoRA on the \textit{current} task. We observe this for \pend{failure_mode_tasks} (see \S\ref{sec:results:ablations}). Practitioners deploying \tglora{} for a task whose knowledge is known to reside in late MLPs (e.g., narrow factual domains) should raise $\tau$ to include more vulnerable components, trading a small forgetting regression for current-task accuracy.

\paragraph{Positioning vs.\ LOFIT and concurrent component-level work.} LOFIT \citep{yin2024lofit} also localises LoRA updates to a subset of components, but it picks them by \textit{new-task adaptation leverage}. \tglora{} picks them by \textit{prior-task resistance}. The two objectives point at overlapping but distinct subsets (Jaccard $=\pend{lofit_tglora_jaccard}$); combining them is a natural follow-up. Concurrent work \citep{anon2026mechforget, lrdistribution2024} has begun to characterise forgetting at component granularity, but via gradient magnitudes and per-layer learning-rate scaling respectively, not via PEFT target selection.
